\documentclass[10pt,twocolumn]{article}

% ── 패키지 ──────────────────────────────────────────────────────────────
\usepackage[margin=1in]{geometry}
\usepackage{amsmath,amssymb,amsthm}
\usepackage{algorithm}
\usepackage{algpseudocode}
\usepackage{booktabs}
\usepackage{graphicx}
\usepackage{xcolor}
\usepackage{hyperref}
\usepackage{microtype}
\usepackage{bm}
\usepackage{enumitem}
\usepackage{caption}
\usepackage{subcaption}

% ── 수식 단축 ────────────────────────────────────────────────────────────
\newcommand{\R}{\mathbb{R}}
\newcommand{\norm}[1]{\left\|#1\right\|}
\newcommand{\cosim}{\mathrm{sim}}

\DeclareMathOperator*{\argmax}{arg\,max}

% ── 제목 ─────────────────────────────────────────────────────────────────
\title{\textbf{Multi-Interest User Modeling via Streaming Text Embeddings\\
for Personalized Ad Recommendation}\\[6pt]
{\normalsize AI506 Data Mining and Search --- Term Project Report}}

\author{\textit{Model-1: Simple Text Embedding Baseline}}
\date{May 2026}

% ═══════════════════════════════════════════════════════════════════════════
\begin{document}
\maketitle

% ── Abstract ─────────────────────────────────────────────────────────────
\begin{abstract}
We present a parameter-free personalized ad recommendation system that
leverages \textbf{pre-trained text embeddings} to represent both search
queries and ads in a shared semantic space, and maintains per-user
\emph{multi-interest vectors} updated online from streaming search and
click events.
Empirical analysis confirms that text-based semantic similarity is
strongly predictive of user preference in ad ranking
(Cohen's $d = 1.56$ between ground-truth and random candidates),
motivating a hybrid scoring function that combines query--ad text similarity
with a multi-faceted user interest profile.
On the ad recommendation validation set (214 queries, 17{,}518 candidates),
our method achieves NDCG@3~$= 0.1228$, a $\mathbf{6{\times}}$ improvement
over the category-wise HistCTR baseline (NDCG@3~$= 0.0211$),
with no trainable parameters.
\end{abstract}

% ═══════════════════════════════════════════════════════════════════════════
\section{Introduction}
\label{sec:intro}

Personalized ad recommendation is a fundamental problem in information
retrieval.
Given a user's search query and behavioral history, the system must rank
thousands of candidate ads such that the most relevant ones appear at the
top.
Two ingredients are central to this problem:
(i) an accurate measure of \emph{relevance} between a query and an ad,
and (ii) a model of the user's \emph{long-term preferences} that
transcends any single query.

Our approach addresses both with a unified text-embedding framework.
Concretely, we make the following contributions:
\begin{enumerate}[leftmargin=*,nosep,itemsep=2pt]
  \item We demonstrate through exploratory data analysis (EDA) that
    \textbf{text embedding similarity} is the strongest signal for user
    preference alignment in this dataset, outperforming category-level
    statistics.
  \item We motivate and instantiate a \textbf{multi-interest user model}
    in text embedding space, where each user is represented by $k$
    interest vectors that collectively capture diverse facets of
    preference --- a design informed by prior work on multi-interest
    recommendation~\cite{li2019mind,cen2020comirec}.
  \item We propose an online soft-assignment update rule and an
    adaptive mixing strategy that gracefully handles cold-start users.
\end{enumerate}

% ═══════════════════════════════════════════════════════════════════════════
\section{Problem Formulation}
\label{sec:problem}

Let $\mathcal{U}$ be the set of users and $\mathcal{A}$ the global ad
corpus of $N = 17{,}518$ ads.
The system observes a \emph{streaming event log} in chronological order;
each event is a tuple
$(u,\, q,\, \{(a_1, y_1), \ldots, (a_m, y_m)\})$,
where $u \in \mathcal{U}$, $q$ is the issued search query,
$a_j \in \mathcal{A}$ are the ads served, and $y_j \in \{0, 1\}$
indicates whether $a_j$ was clicked.

\paragraph{Task B --- Ad Recommendation (NDCG@3).}
Given a held-out event $(u, q)$,
rank all $N$ candidate ads and return a top-$K$ list.
Performance is measured by Normalized Discounted Cumulative Gain:
\[
  \mathrm{NDCG@3} = \frac{1}{\log_2(r+1)} \cdot \mathbf{1}[r \leq 3],
\]
where $r$ is the rank of the single ground-truth ad.
Unlike click prediction (Task A), this task is a \emph{ranking} problem:
the goal is to surface the most \emph{preferred} ad for the user,
which requires modeling long-term user preference beyond the current query.

% ═══════════════════════════════════════════════════════════════════════════
\section{Motivation}
\label{sec:motivation}

\subsection{Text Embeddings as a Shared Semantic Space}
\label{sec:motivation:emb}

Both search queries and ad titles are natural-language text.
This seemingly simple observation has an important consequence:
when the same pre-trained language model encodes both queries and ads,
they are projected into the \emph{same semantic space} with no
cross-modal gap.

Formally, we use a single encoder $\phi : \text{text} \to \R^d$
($d = 384$, all-MiniLM-L6-v2 family) and represent every entity as:
\[
  \phi(q) \in \R^d, \quad \phi(a) \in \R^d.
\]
Because $q$ and $a$ are both text, their embeddings are directly
comparable via cosine similarity without any additional projection layer
or domain adaptation.
This shared representation enables the model to measure the
\textbf{semantic alignment} between what a user is searching for and
what an ad is offering, using a single inner product.

\subsection{EDA: Text Similarity Captures User Preference}
\label{sec:motivation:eda}

To empirically validate the choice of text embeddings, we conduct an
exploratory analysis on the ad recommendation validation set.
For each of the 214 validation queries, we compute the cosine similarity
$\cosim(\hat{\phi}(q), \hat{\phi}(a))$ for (i) the ground-truth ad
and (ii) five randomly sampled candidate ads from the pool of 17{,}518.

\begin{table}[h]
\centering
\caption{Discriminability of text similarity for Task B.
         Cohen's $d$ measures how well text similarity separates
         the preferred ad from random candidates.}
\label{tab:eda}
\small
\begin{tabular}{lccc}
\toprule
Group & $n$ & Mean sim. & Std \\
\midrule
Ground-truth (preferred) ads & 214 & \textbf{0.548} & 0.213 \\
Randomly sampled ads         & 1{,}070 & 0.253 & 0.160 \\
\midrule
\multicolumn{2}{l}{Gap (preferred $-$ random)} & \multicolumn{2}{c}{$+0.295$} \\
\multicolumn{2}{l}{Cohen's $d$} & \multicolumn{2}{c}{$\mathbf{1.56}$ (large effect)} \\
\bottomrule
\end{tabular}
\end{table}

As shown in Table~\ref{tab:eda}, preferred ads have a mean similarity
of $0.548$ to their query, while randomly sampled ads score only $0.253$
--- a gap of $+0.295$ with an effect size of Cohen's $d = 1.56$.
By standard convention, $d > 0.8$ indicates a large effect;
$d = 1.56$ thus demonstrates that text embedding similarity is a
\textbf{highly discriminative signal} for user preference in this domain.

This result underpins our design decision to use text embeddings as the
primary representation: no other feature in the dataset (e.g., category
tags, historical CTR) provides as direct a measure of relevance between
a user's expressed intent and an ad's content.

\subsection{Ad Recommendation as a User Preference Ranking Problem}
\label{sec:motivation:ranking}

Task B is fundamentally a \textbf{preference ranking} problem.
Among 17{,}518} candidate ads, the model must identify the one that best
matches the user's immediate intent \emph{and} long-term interests.
The current query captures immediate intent, but it alone cannot account
for user-specific preference patterns accumulated over time.

For instance, two users issuing the identical query ``running shoes''
may prefer very different ad styles: one favors minimalist trail shoes,
the other high-tech road-running gear.
A query-only model ranks both users identically, whereas a model that
maintains a personalized user profile can distinguish them.

This motivates the inclusion of a \emph{user interest component} in the
scoring function, operating alongside the text similarity term.

% ═══════════════════════════════════════════════════════════════════════════
\section{Related Work: Multi-Interest Recommendation}
\label{sec:related}

A single user-embedding vector is insufficient to represent the
diversity of real-world user interests~\cite{li2019mind}.
A user who purchases both sports equipment and cooking supplies cannot
be accurately described by a single point in the embedding space.

The multi-interest paradigm addresses this by maintaining \emph{multiple}
prototype vectors per user.
\textbf{MIND}~\cite{li2019mind} introduced dynamic routing to assign
items to different interest capsules.
\textbf{ComiRec}~\cite{cen2020comirec} proposed a controllable
multi-interest framework that disentangles user representations for
recommendation diversity.
These works consistently show that multi-vector user representations
outperform single-vector baselines in retrieval recall and ranking quality.

Our work adapts this paradigm to the text-embedding domain.
Unlike prior methods that operate on ID-based item embeddings learned
end-to-end, we maintain interest vectors in the \emph{same pre-trained
text embedding space} used for queries and ads.
This design choice means the interest vectors are directly comparable
to both queries and ad embeddings via cosine similarity --- no learned
projection or compatibility layer is required.
The online soft-assignment update rule (Section~\ref{sec:update})
further distinguishes our approach by enabling streaming updates without
batch retraining.

% ═══════════════════════════════════════════════════════════════════════════
\section{Method}
\label{sec:method}

Our model, \textbf{MultiInterest}, integrates the insights above into
three components:
(i) a multi-vector user interest representation in text embedding space,
(ii) an online soft-assignment update rule,
and (iii) a hybrid scoring function combining text similarity and
personalized interest.

% ─────────────────────────────────────────────────────────────────────────
\subsection{Multi-Interest User Representation}
\label{sec:repr}

Each user $u$ is associated with $k$ \emph{interest vectors}
$\bm{V}^u = [\bm{v}^u_1, \ldots, \bm{v}^u_k] \in \R^{k \times d}$,
residing in the same $d$-dimensional space as query and ad embeddings.

A single vector collapses all of a user's interests into one direction,
which is inadequate when users exhibit multiple distinct preference
facets (Section~\ref{sec:motivation:ranking}).
With $k > 1$ vectors, each prototype can specialize to a distinct
facet of the user's interest profile, as established in the
multi-interest recommendation literature~\cite{li2019mind,cen2020comirec}.

\paragraph{Initialization.}
When user $u$ is first encountered, each vector is drawn i.i.d.\ from
a standard Gaussian and $\ell_2$-normalized:
\[
  \bm{v}^u_i \leftarrow \frac{\bm{z}_i}{\norm{\bm{z}_i}}, \quad
  \bm{z}_i \sim \mathcal{N}(\bm{0}, \bm{I}_d).
\]
Random initialization ensures the $k$ vectors initially point in
different directions, promoting specialization as behavioral evidence
accumulates.
Zero initialization is avoided because it makes all soft-assignment
weights uniform, causing all $k$ vectors to collapse to the same direction.

% ─────────────────────────────────────────────────────────────────────────
\subsection{Online Soft-Assignment Interest Update}
\label{sec:update}

When a new text embedding $\bm{e}$ arrives (from a search query or
clicked ad), each interest vector $\bm{v}^u_i$ is updated proportionally
to its \emph{soft responsibility} for $\bm{e}$.
All operations take place in the shared text embedding space.

\paragraph{Step 1 --- Soft assignment weights.}
The cosine distance from each interest vector to the normalized embedding
$\hat{\bm{e}} = \bm{e}/\norm{\bm{e}}$ is:
\[
  d_i = 1 - \hat{\bm{v}}^u_i \cdot \hat{\bm{e}},
  \quad \hat{\bm{v}}^u_i = \bm{v}^u_i / \norm{\bm{v}^u_i}.
\]
The distances are converted to soft weights via a temperature-scaled
exponential:
\begin{equation}
  w_i = \frac{\exp(-d_i / \tau)}{\displaystyle\sum_{j=1}^{k} \exp(-d_j / \tau)},
  \label{eq:softweight}
\end{equation}
where $\tau > 0$ controls the sharpness.
As $\tau \to 0$ the assignment becomes hard (nearest neighbor);
as $\tau \to \infty$ it becomes uniform.

\paragraph{Step 2 --- Weighted additive update.}
Each interest vector moves toward (or away from) $\hat{\bm{e}}$,
weighted by responsibility and scaled by event-specific learning rate $\alpha$:
\begin{equation}
  \bm{v}^u_i \leftarrow
    \bm{v}^u_i + \sigma \cdot \alpha \cdot w_i \cdot \hat{\bm{e}},
  \label{eq:update}
\end{equation}
where $\sigma \in \{+1, -1\}$ indicates attraction or repulsion.
Table~\ref{tab:update_rules} summarizes the update rules.

\begin{table}[h]
\centering
\caption{Interest update rules per event type.
         All embeddings belong to the same pre-trained text embedding space.}
\label{tab:update_rules}
\small
\begin{tabular}{lccc}
\toprule
Event & $\sigma$ & $\alpha$ & Intuition \\
\midrule
Search query $q$ & $+1$ & $\alpha_{\text{search}}$ & Drift toward topic \\
Ad click $a$     & $+1$ & $\alpha_{\text{click}}$  & Strengthen preference \\
Ad non-click $a$ & $-1$ & $\alpha_{\text{neg}}$    & Suppress dispreference \\
\bottomrule
\end{tabular}
\end{table}

The full procedure is in Algorithm~\ref{alg:update}.

\begin{algorithm}[t]
\caption{Online Interest Update}
\label{alg:update}
\begin{algorithmic}[1]
\Require User $u$, text embedding $\bm{e}$, sign $\sigma$, rate $\alpha$, temp.\ $\tau$
\State $\hat{\bm{e}} \leftarrow \bm{e} / \norm{\bm{e}}$
  \Comment{normalize in shared embedding space}
\State Retrieve (or initialize) $\bm{V}^u = [\bm{v}^u_1,\ldots,\bm{v}^u_k]$
\For{$i = 1$ \textbf{to} $k$}
  \State $d_i \leftarrow 1 - (\bm{v}^u_i / \norm{\bm{v}^u_i}) \cdot \hat{\bm{e}}$
  \State $r_i \leftarrow \exp(-d_i / \tau)$
\EndFor
\State $\bm{w} \leftarrow \bm{r} / \sum_j r_j$ \Comment{Eq.~\eqref{eq:softweight}}
\For{$i = 1$ \textbf{to} $k$}
  \State $\bm{v}^u_i \leftarrow \bm{v}^u_i
    + \sigma \cdot \alpha \cdot w_i \cdot \hat{\bm{e}}$
    \Comment{Eq.~\eqref{eq:update}}
\EndFor
\end{algorithmic}
\end{algorithm}

% ─────────────────────────────────────────────────────────────────────────
\subsection{Hybrid Scoring for Ad Recommendation}
\label{sec:scoring}

Given a test event $(u, q)$ and the full candidate set $\mathcal{A}$,
we assign a relevance score to each candidate ad $a_j$:
\begin{equation}
  s(u, q, a_j) =
    \underbrace{(1-\gamma)\,\cosim\!\left(\hat{\phi}(q),\, \hat{\phi}(a_j)\right)}_{\text{immediate intent}}
    +\; \underbrace{\gamma \max_{i=1}^{k}\,
    \cosim\!\left(\hat{\bm{v}}^u_i,\, \hat{\phi}(a_j)\right)}_{\text{long-term preference}},
  \label{eq:score}
\end{equation}
where $\hat{\phi}(\cdot)$ denotes the $\ell_2$-normalized text embedding
and $\gamma \in [0,1]$ controls the balance between the two terms.

\paragraph{First term --- Query--ad text similarity (immediate intent).}
Because queries and ads are both encoded by the same text encoder,
$\cosim(\hat{\phi}(q), \hat{\phi}(a_j))$ measures semantic alignment
directly in the text embedding space, \emph{without any learned
projection}.
The EDA result (Table~\ref{tab:eda}, Cohen's $d = 1.56$) confirms
this is the dominant signal for user preference in ad ranking.

\paragraph{Second term --- Interest--ad similarity (long-term preference).}
$\max_{i}\,\cosim(\hat{\bm{v}}^u_i, \hat{\phi}(a_j))$ personalizes the
ranking by finding the most relevant interest facet for each candidate.
Because the interest vectors live in the same text embedding space,
this comparison is also computed as a plain cosine similarity ---
no additional compatibility layer is needed.
The $\max$ aggregation is preferred over $\text{mean}$ because
(a) it selects the most applicable interest facet rather than
averaging over unrelated ones, and
(b) it is more robust when some interest vectors are still near
their random initialization.

\paragraph{Matrix form.}
All $N$ candidates are scored simultaneously via:
\[
  \bm{s}_u = (1-\gamma)\,\bm{C}\hat{\phi}(q)
           + \gamma \max_{i}\,\bigl(\bm{C}\hat{\bm{V}}^{u\top}\bigr)_{\cdot i},
\]
where $\bm{C} \in \R^{N \times d}$ is the pre-computed matrix of
normalized ad text embeddings.
The $\max$ is taken column-wise over the $k$ interest similarities.

% ─────────────────────────────────────────────────────────────────────────
\subsection{Adaptive Gamma for Cold-Start Users}
\label{sec:adaptive}

In the training corpus, only 14\% of users ($2{,}385$ of $16{,}975$)
have at least one click event.
For users without click history, the interest vectors $\bm{V}^u$ remain
near their random initialization and carry no meaningful preference
signal.
We therefore adopt an \emph{adaptive} mixing coefficient:
\begin{equation}
  \gamma_{\text{eff}}(u) =
  \begin{cases}
    \gamma & \text{if user } u \text{ has } \geq\!1 \text{ click in training,} \\
    0      & \text{otherwise (cold-start).}
  \end{cases}
  \label{eq:adaptive}
\end{equation}
For cold-start users the score reduces to pure query--ad text similarity,
which the EDA identifies as the strongest available signal
(Table~\ref{tab:eda}).

% ─────────────────────────────────────────────────────────────────────────
\subsection{Search-to-Click Learning Rate Ratio}
\label{sec:design}

A subtle but important design choice is the relative learning rate
assigned to search and click events.
In the training log there are $276{,}807$ search events but only
$3{,}560$ click events --- a $78:1$ imbalance.
With equal learning rates the accumulated search updates would dominate,
biasing interest vectors toward generic query topics rather than
specific click preferences.

We set $\alpha_{\text{search}} = 0.01$ and $\alpha_{\text{click}} = 0.5$
(a 50-fold difference), so each click contributes substantially more to
the interest vector than each individual search event.
Empirically, reducing $\alpha_{\text{search}}$ from $0.1$ to $0.01$
improves NDCG@3 from $0.1065$ to $0.1228$ (+15\%, Table~\ref{tab:ablation}).

\begin{table}[h]
\centering
\caption{Ablation on $\alpha_{\text{search}}$ (ad recommendation, NDCG@3).}
\label{tab:ablation}
\small
\begin{tabular}{lcc}
\toprule
Configuration & AUC$_{\text{rank}}$ & NDCG@3 \\
\midrule
$\alpha_{\text{search}} = 0.1$ & 0.8417 & 0.1065 \\
$\alpha_{\text{search}} = 0.01$ \textbf{(ours)} & \textbf{0.8500} & \textbf{0.1228} \\
\bottomrule
\end{tabular}
\end{table}

% ─────────────────────────────────────────────────────────────────────────
\subsection{Hyperparameter Summary}
\label{sec:hyper}

\begin{table}[h]
\centering
\caption{Final hyperparameter configuration.}
\label{tab:hyper}
\small
\begin{tabular}{lcc}
\toprule
Parameter & Symbol & Value \\
\midrule
\# interest vectors    & $k$                      & 5 \\
Embedding dimension    & $d$                      & 384 \\
Search update rate     & $\alpha_{\text{search}}$ & 0.01 \\
Click update rate      & $\alpha_{\text{click}}$  & 0.50 \\
Neg.\ update rate      & $\alpha_{\text{neg}}$    & 0.00 \\
Temperature            & $\tau$                   & 1.0 \\
Interest mixing weight & $\gamma$                 & 0.5 (adaptive) \\
\bottomrule
\end{tabular}
\end{table}

% ═══════════════════════════════════════════════════════════════════════════
\section{Complexity}
\label{sec:complexity}

\paragraph{Training.}
Each event costs $O(kd)$ (soft-weight computation and update).
Total training cost: $O(|\mathcal{E}| \cdot kd)$ with no gradient
computation or matrix inversion.

\paragraph{Inference.}
Scoring all $N$ candidates requires $O(Nkd)$.
The ad embedding matrix $\bm{C}$ is pre-computed once and reused.
For our setting ($N\!=\!17{,}518$, $k\!=\!5$, $d\!=\!384$),
this is ${\approx}\,3.4 \times 10^7$ operations,
achievable in under 1 second on a standard CPU.

% ═══════════════════════════════════════════════════════════════════════════
\section{Experiments}
\label{sec:results}

\paragraph{Setup.}
We train on $320{,}000$ (search, ad, click) triples in chronological order.
Evaluation uses \texttt{ad\_validation}: 214 queries, one ground-truth
ad each, ranked against all $17{,}518$ candidates.

\paragraph{Baselines.}
\begin{itemize}[leftmargin=*,nosep,itemsep=1pt]
  \item \textbf{Random}: uniform random score.
  \item \textbf{HistCTR}: rank by historical CTR within category
    (official baseline).
  \item \textbf{Query-only} ($\gamma\!=\!0$): pure query--ad text
    similarity, no user interest.
\end{itemize}

\begin{table}[h]
\centering
\caption{Ad recommendation results (NDCG@3).}
\label{tab:results}
\small
\begin{tabular}{lcc}
\toprule
Method & NDCG@3 & Rank-1 hits \\
\midrule
Random & 0.0008 & 0 \\
HistCTR (baseline) & 0.0211 & 2 \\
Query-only ($\gamma\!=\!0$) & 0.0989 & 14 \\
\textbf{MultiInterest (ours)} & \textbf{0.1228} & \textbf{17} \\
\bottomrule
\end{tabular}
\end{table}

\begin{table}[h]
\centering
\caption{Rank distribution of ground-truth ads ($n\!=\!214$).}
\label{tab:rankdist}
\small
\begin{tabular}{lrc}
\toprule
Rank & Count & NDCG weight \\
\midrule
1       & 17  & 1.000 \\
2       & 6   & 0.631 \\
3       & 2   & 0.500 \\
$> 3$ & 189 & 0.000 \\
\bottomrule
\end{tabular}
\end{table}

The query-only baseline already surpasses HistCTR by $4.7\times$,
confirming that \textbf{text embedding similarity is the dominant
relevance signal} for this task.
Adding the multi-interest term yields a further $+24\%$ relative
improvement ($0.0989 \to 0.1228$), demonstrating that personalized
user preference provides complementary signal beyond query intent alone.

% ═══════════════════════════════════════════════════════════════════════════
\section{Conclusion}
\label{sec:conclusion}

We presented \textbf{MultiInterest}, a streaming multi-interest model
for personalized ad recommendation built entirely on pre-trained
\textbf{text embeddings}.
Three properties make text embeddings particularly suitable for this task:
(1) queries and ads share the same natural-language modality, enabling
direct cosine similarity without cross-modal projection;
(2) EDA confirms that text-based similarity is highly predictive of
user preference (Cohen's $d = 1.56$);
and (3) interest vectors maintained in the same embedding space are
immediately comparable to both query and ad representations.

The multi-interest architecture, grounded in prior work on multi-vector
user modeling~\cite{li2019mind,cen2020comirec}, captures diverse facets
of user preference through $k$ interest prototypes, updated online via
a soft-assignment rule.
Combined with adaptive gamma and a carefully calibrated
search-to-click learning rate ratio, our method achieves
NDCG@3~$= 0.1228$ --- a $\mathbf{6\times}$ improvement over the
HistCTR baseline --- with no trainable parameters.

% ── 참고문헌 ──────────────────────────────────────────────────────────────
\begin{thebibliography}{9}

\bibitem{li2019mind}
  Chao Li, Zhiyuan Liu, Mengmeng Wu, Yuchi Xu, Huan Zhao, Pipei Huang,
  Guoliang Kang, Qiwei Chen, Wei Li, and Dik Lun Lee.
  \newblock Multi-Interest Network with Dynamic Routing for Recommendation
    at Tmall.
  \newblock In \textit{Proceedings of the 28th ACM International Conference
    on Information and Knowledge Management (CIKM)}, 2019.

\bibitem{cen2020comirec}
  Yukuo Cen, Jianwei Zhang, Xu Zou, Chang Zhou, Hongxia Yang, and Jie Tang.
  \newblock Controllable Multi-Interest Framework for Recommendation.
  \newblock In \textit{Proceedings of the 26th ACM SIGKDD International
    Conference on Knowledge Discovery \& Data Mining}, 2020.

\end{thebibliography}

\end{document}
